\section[p-Brauer character]{$\boldsymbol{p}$-Brauer character}\label{section5}

In this section, we will define a new notion of generalized Brauer character and prove properties of it. We will give a counterexample to a conjecture of Borcherds about vanishing of Tate cohomology.

Let $p$ be a prime number. Let $R$ be a discrete valuation ring of mixed characteristic $(0,p)$, $v\colon R\rightarrow \mathbb{Z}_{\geq 0}\cup \{ \infty \}$ be a~surjective valuation on $R$ and $K=\operatorname{Frac}(R)$.

We can't use the Brauer character for $R_p/(1-\zeta_N^n)R_p$ because this is not a vector space over a field. To solve this, we will define a $p$-Brauer character.

\begin{Definition}\label{def5.1}
Let $G$ be a finite group and $M$ be a finite length $R[G]$-module. Let $h$ be a $p$-regular element of $G$. Suppose that $M$ has a composition series $0=M_0\subsetneq M_1\subsetneq \dots \subsetneq M_n=M$ as an $R[h]$-module.
 We define the $p$-Brauer character of $M$ for $h$ by
\[
\widetilde{\operatorname{Tr}}_p(h|M)=\cfrac{1}{v(p)} \sum_{i=1}^{n}\operatorname{Tr} \big(h|\widetilde{M_i}\otimes K\big),
\]
 where $\widetilde{M_i}$ is an $R$-free $R[h]$-module satisfying $M_i/M_{i-1} \cong\widetilde{{M}_i}\otimes_{R[h]}R/\mathfrak{m}_R[h]$.
\end{Definition}

\begin{Lemma}
This definition does not depend on our choices of $M_i$ and $\widetilde{M}_i$.
\end{Lemma}

\begin{proof}
For $M_i/M_{i-1} \cong\widetilde{{M}_i}\otimes_{R[h]}R/\mathfrak{m}_R[h]\cong\widetilde{{M}_i'}\otimes_{R[h]}R/\mathfrak{m}_R[h]$, we have $\operatorname{Tr}\big(h|\widetilde{M}_i\otimes K\big)=\widetilde{\operatorname{Tr}}(h|M_i/M_{i-1})=\operatorname{Tr}\big(h|\widetilde{M}_i'\otimes K\big)$. By the Jordan--H\"older theorem, the composition series of~$M$ is unique up to permutation and isomorphism. That is, for any two composition series $0=M_0\subsetneq M_1\subsetneq \dots \subsetneq M_n=M$, $0=M_0'\subsetneq M_1'\subsetneq \dots \subsetneq M_{n'}'=M$, we have $n=n'$ and some permutation $\sigma$ of $\{1,\dots,n\}$ s.t.\ $M_i/M_{i-1}\cong M_{\sigma(i)}'/M_{\sigma(i)-1}'$. Hence, $\sum_{i=1}^{n}\widetilde{\operatorname{Tr}}(h|M_i/M_{i-1})=\sum_{i=1}^{n}\widetilde{\operatorname{Tr}}(h|M_{\sigma(i)}'/M_{\sigma(i)-1}') =\sum_{i=1}^n\widetilde{\operatorname{Tr}}(h|M_i'/M_{i-1}')$. Therefore the $p$-Brauer character is well-defined.
\end{proof}

Note that it does depend on a choice of a prime number~$p$.

We can consider the notion of an irreducible $p$-Brauer character, so that $M$ is a simple $R[h]$-module that $M\otimes_{R[h]} R/\mathfrak{m}_R[h]$ is a simple $R/\mathfrak{m}_R[h]$-module. In particular, an irreducible $p$-Brauer is an irreducible Brauer character and a character table of $p$-Brauer characters is the same as of Brauer characters.

We will prove that the $p$-Brauer character is additive, where ``additive'' means that for any short exact sequence $0\rightarrow M\rightarrow M'\rightarrow M''\rightarrow 0$, $\widetilde{\operatorname{Tr}}_p$ satisfies $\widetilde{\operatorname{Tr}}_p(h|M')=\widetilde{\operatorname{Tr}}_p(h|M)+\widetilde{\operatorname{Tr}}_p(h|M'')$.

\begin{Lemma}\label{additive}
The $p$-Brauer character is additive on short exact sequences.
\end{Lemma}

\begin{proof}
Let $0\rightarrow M\xrightarrow{\varphi} M'\xrightarrow{\psi} M''\rightarrow 0$ be an exact sequence of finite length $R[h]$-modules. Suppose that $M$ has length $m$ and a composition series $0=M_0\subsetneq M_1\subsetneq \dots \subsetneq M_m=M$. We use analogous notation for~$M'$ and~$M''$. Then the length~$m'$ of~$M'$ is $m+m''$ and
\[
0=\varphi(M_0)\subsetneq \varphi(M_1)\subsetneq \dots \subsetneq \varphi(M_m)=\psi^{-1}(M''_0)\subsetneq \psi^{-1}(M''_1)\subsetneq \dots \subsetneq \psi^{-1}(M''_{m''})=M'
\]
 is a composition series of $M'$. By injectivity of $\varphi$, $M_i\cong \varphi(M_i)$ for any $i \in \{0, \dots,m\}$ and by surjectivity of $\psi$, $\psi^{-1}(M''_j)/\operatorname{Ker}(\psi)\cong M''_j$ for any $j \in \{0, \dots,m''\}$. Hence, we have $\varphi(M_i)/\varphi(M_{i-1})\cong M_i/M_{i-1}$ and
\[
\psi^{-1}(M''_j)/\psi^{-1}(M''_{j-1})\cong \big(\psi^{-1}(M''_j)/\operatorname{Ker}(\psi)\big)/\big(\psi^{-1}(M''_{j-1})/\operatorname{Ker}(\psi)\big) \cong M''_j/M''_{j-1}.
\]
 Therefore,
\begin{align*}
\widetilde{\operatorname{Tr}}_p(h|M')&=\frac{1}{v(p)}\sum_{i=1}^m\widetilde{\operatorname{Tr}}\big(h|\varphi(M_i)/\varphi(M_{i-1})\big) +\frac{1}{v(p)}\sum_{j=1}^{m''}\widetilde{\operatorname{Tr}}\big(h|\psi^{-1}(M''_j)/\psi^{-1}(M''_{j-1})\big) \\
 &=\frac{1}{v(p)}\sum_{i=1}^m\widetilde{\operatorname{Tr}}(h|M_i/M_{i-1})+\frac{1}{v(p)}\sum_{j=1}^{m''}\widetilde{\operatorname{Tr}}(h|M''_j/M''_{j-1}) \\
 &=\widetilde{\operatorname{Tr}}_p(h|M)+\widetilde{\operatorname{Tr}}_p(h|M'').\tag*{\qed}
\end{align*}\renewcommand{\qed}{}
\end{proof}

\begin{Lemma}\label{tensorBch}
Let $S$ be the ring of integers of a finite extension of $\operatorname{Frac}(R)$ and $A$ be a finite length $R[h]$-module. Then $\widetilde{\operatorname{Tr}}_p(h|A\otimes_{R}S)=\widetilde{\operatorname{Tr}}_p(h|A)$.
\end{Lemma}

\begin{proof}
We have that the $p$-Brauer character of $A$ is described the sum of $p$-Brauer characters of composition factors of $A$ by Definition~\ref{def5.1}. Hence, it suffices to consider the case $A$ is a simple $R[h]$-module on which $h$ acts as $\zeta_{|h|}^m$. Now,
\[
\widetilde{\operatorname{Tr}}_p(h|A)=\frac{1}{v_R(p)}\operatorname{Tr} \big(h|\widetilde{A}\otimes \operatorname{Frac}(R)\big),
\]
where $v_R$ is a discrete valuation on~$R$, $\widetilde{A}$ is an $R$-free $R[h]$-module satisfying $A\cong \widetilde{A}\otimes_{R[h]} R/\mathfrak{m}_R[h]$ and $\mathfrak{m}_R$ is a maximal ideal of~$R$. We have
\[
A\otimes_R S\cong \widetilde{A}\otimes_{R[h]} R/\mathfrak{m}_R[h]\otimes_R S\cong \widetilde{A}\otimes_{R[h]} S/\mathfrak{m}_RS[h].
\]

 Let $v_S$ be a discrete valuation on $S$ and $\mathfrak{m}_S$ be a maximal ideal of~$S$. If $S$ is an unramified extension, then $\mathfrak{m}_S=\mathfrak{m}_RS$ and $v_S(p)=v_R(p)$. Hence,
 \begin{align*}
\widetilde{\operatorname{Tr}}_p(h|A\otimes_R S) &=\frac{1}{v_S(p)}\operatorname{Tr}\big(h|\widetilde{A}\otimes_R S\otimes \operatorname{Frac}(S)\big)\\
& =\frac{1}{v_R(p)}\operatorname{Tr}\big(h|\widetilde{A}\otimes \operatorname{Frac}(R)\big)=\widetilde{\operatorname{Tr}}_p(h|A).
 \end{align*}

 If $S$ is a ramified extension, then there is a positive integer $k>1$ s.t.\ $\mathfrak{m}_S^k=\mathfrak{m}_RS$ and $v_S(p)=kv_R(p)$. Hence, we have $A\otimes_R S=\widetilde{A}\otimes_{R[h]}S/\mathfrak{m}_S^k[h]$ and a composition series $0\subsetneq \widetilde{A}\otimes S/\mathfrak{m}_S\subsetneq \dots \subsetneq \widetilde{A}\otimes S/\mathfrak{m}_S^k$. Its composition factors are all isomorphic to $\widetilde{A}\otimes_{R[h]} S/\mathfrak{m}_S[h]$. Therefore,
\begin{align*}
\widetilde{\operatorname{Tr}}_p(h|A\otimes_R S) & =\frac{1}{v_S(p)}\sum_{i=1}^k\operatorname{Tr} \big(h|\widetilde{A}\otimes_R S\otimes \operatorname{Frac}(S)\big)\\
& =\frac{1}{kv_R(p)}k\operatorname{Tr}\big(h|\widetilde{A}\otimes \operatorname{Frac}(R)\big) =\widetilde{\operatorname{Tr}}_p(h|A).\tag*{\qed}
\end{align*}\renewcommand{\qed}{}
\end{proof}

By Proposition \ref{tensor} and Lemma~\ref{tensorBch}, any finitely generated $\mathbb{Z}$-free $\mathbb{Z}[\langle g,h \rangle]$-module base changes to an $R_p[\langle g,h \rangle]$-module that is $R_p$-free of finite rank and the $p$-Brauer character of the Tate cohomology is unchanged.
By Lemma \ref{additive}, it suffices to consider the $p$-Brauer character for the Tate cohomology of~$R_{n,m}$.

We will calculate the Tate cohomology of $R_{n,m}$ by using the $p$-Brauer character. To do this, we use the following lemma:

\begin{Lemma}[{\cite[Chapter III, Lemma 3]{CF}}] \label{pR}
 Let $q=p^k$ for a positive integer~$k$. Then $pR=(1-\zeta_q)^{\varphi(q)}R$, where $\varphi$ is Euler's totient function.
\end{Lemma}

For a positive integer $s$ and $R=\mathbb{Z}_{p}[\zeta_{p^s}]$, a discrete valuation $v(p)$ is given by $\varphi(p^s)$ and maximal ideal of~$R$ is $(1-\zeta_{p^s})R$.

\begin{Lemma}\label{Bch}
Suppose that the prime factorization of $N$ is $\prod p_i^{n_i}$.
Then,
\[
\widetilde{\operatorname{Tr}}_{p_i}(h|\hat{H}^*(g,R_{n,m}))= \begin{cases}
n_i\zeta_{|h|}^m, & n=0, \vspace{1mm}\\
\dfrac{-\zeta_{|h|}^m}{\varphi \big(p_i^l\big)}, & |g^n|=p_i^l, \ 1\leq l\leq n_i, \vspace{1mm}\\
0, & otherwise.
\end{cases}.
\]
\end{Lemma}

\begin{proof}Now, $h$ acts as $\zeta_{|h|}^m$ on $\hat{H}^*(g,R_{n,m})$ and $v(p)=\varphi\big(p_i^{n_i}\big)$. By \[
\widetilde{\operatorname{Tr}}_{p_i}\big(h|\hat{H}^*(g,R_{n,m})\big)=\widetilde{\operatorname{Tr}}_{p_i}\big(h|\hat{H}^0(g,R_{n,m})\big) -\widetilde{\operatorname{Tr}}_{p_i}\big(h|\hat{H}^1(g,R_{n,m})\big)
\]
 and Lemma~\ref{Tcoho}, it suffices to calculate the $p_i$-Brauer characters of $R_{p_i}/NR_{p_i}$ and $R_{p_i}/(1-\zeta_N^n)R_{p_i}$.

 When $n=0$, by Lemmas~\ref{Tcoho} and~\ref{pR},
\[
 R_{p_i}/NR_{p_i}\cong R_{p_i}/p_i^{n_i}R_{p_i}\cong R_{p_i}/\big(1-\zeta_{p_i^{n_i}}\big)^{n_i\varphi(p_i^{n_i})}R_{p_i}.
\]
This module has a composition series
\[
0\subsetneq R_{p_i}/(1-\zeta_{{p_i}^{n_i}})R_{p_i}\subsetneq \dots \subsetneq R_{p_i}/\big(1-\zeta_{{p_i}^{n_i}}\big)^{n_i\varphi(p_i^{n_i})}R_{p_i}
\]
 and its composition factors are all isomorphic to $R_{p_i}/\big(1-\zeta_{{p_i}^{n_i}}\big)R_{p_i}$. Hence,
\[
 \widetilde{\operatorname{Tr}}_{p_i}(h|R_{p_i}/NR_{p_i})=n_i\zeta_{|h|}^m.
\]

 When $|g^n|=p_i^l$, $\zeta_N^n$ is a $p_i^l$-th primitive root of unity. By Lemma~\ref{pR}, we have
\[
 \big(1-\zeta_{p_i^l}\big)^{\varphi(p_i^l)}=\big(1-\zeta_{p_i^{n_i}}\big)^{\varphi(p_i^{n_i})}.
\]
 Hence, by Lemma~\ref{Tcoho}, we have
\[
R_{p_i}/\big(1-\zeta_{{p_i}^l}\big)R_{p_i}\cong R_{p_i}/\big(1-\zeta_{{p_i}^{n_i}}\big)^{\varphi(p_i^{n_i})/\varphi(p_i^l)}R_{p_i}.
\]
 This module has the length $\varphi\big(p_i^{n_i}\big)/\varphi\big(p_i^l\big)$ and its composition factors are all isomorphic to $R_{p_i}/\big(1-\zeta_{{p_i}^{n_i}}\big)R_{p_i}$. Calculating the $p_i$-Brauer character, we have
\[
 \widetilde{\operatorname{Tr}}_{p_i}\big(h|R_{p_i}/\big(1-\zeta_{p_i^l}\big)R_{p_i}\big)=1/\varphi\big(p_i^l\big).
\]

 When $|g^n|=p_i^kq$, where $q\neq 1$, $(q,p_i)=1$ and $0\leq k\leq n_i$, $\zeta_N^n$ is a $p_i^kq$-th primitive root of unity. Let $\pi$ be a surjective homomorphism $R_{p_i}\rightarrow \mathbb{F}_{p_i}\big[\zeta_{N|h|/p_i^{n_i}}\big]$; $\sum_{j=0}^\infty c_jp_i^j \mapsto c_0$, $\zeta_{N|h|} \mapsto \zeta_{N|h|/p_i^{n_i}}$. Then $\operatorname{Ker}(\pi)$ is the maximal ideal of $R_{p_i}$ and $\operatorname{Frob}_{p_i}$ on $\mathbb{F}_{p_i}[\zeta_{N|h|/p_i^{n_i}}]$ is an automorphism. We have $\pi\big(\zeta_{p_i^kq}\big)=\operatorname{Frob}_{p_i}^{-k}\big(\zeta_{p_i^kq}^{p_i^k}\big)\neq 1$.
 Hence, $1-\zeta_N^n \notin \operatorname{Ker}(\pi)$. $1-\zeta_N^n$ is a unit of $R_{p_i}$ because every element of a local ring not included in a maximal ideal is a unit. Therefore, $R_p/\big(1-\zeta_N^n\big)R_p=0$.
\end{proof}

We generalize Proposition~\ref{prop2.2} \cite[Proposition~2.2]{BR} using the $p$-Brauer character.

\begin{Theorem} \label{2.2kai}
Let $G$ be a finite group. Suppose that $g \in \operatorname{Cent}(G)$ has order $N$ and that $h \in G$ is an $N$-regular element. Let~$A$ be a finitely generated $R$-free $R[\langle g,h \rangle]$-module. Then $\hat{H}^*(g,A)=\hat{H}^0(g,A)-\hat{H}^1(g,A)$ is a virtual representation of $\langle h \rangle$, and
\begin{gather*}
\widetilde{\operatorname{Tr}}_{p_i}(h|\hat{H}^*(g,A))=\sum_{k=1}^{N-1}a_{k,p_i}\operatorname{Tr}\big(g^kh|A\big) ,
\\
a_{k,p_i}=\begin{cases}
\dfrac{1}{\sum_{(p_i,d)=1,d|N}\varphi (N/d)}\left(n_i-l-\dfrac{n_i-l-1}{p_i}\right), & \big(k,p_i^{n_i}\big)=p_i^l, \ 0\leq l \leq n_i-1, \\
0, & \big(k,p_i^{n_i}\big)=p_i^{n_i},
\end{cases}
\end{gather*}
where the prime factorization of $N$ is $\prod p_i^{n_i}$.
\end{Theorem}

\begin{proof}
Suppose
\[
\widetilde{\operatorname{Tr}}_{p_i}\big(h|\hat{H}^*(g,A)\big)=\sum_{k=1}^{N-1}a_{k,p_i}\operatorname{Tr}\big(g^kh|A\big). \]
It suffices to take $A=R_{n,m}$.
By an inverse discrete Fourier transform and Lemma~\ref{Bch}, we have
\begin{align*}
a_{k,p_i}&=\frac{1}{N}\sum_{b=0}^{N-1}\widetilde{\operatorname{Tr}}_{p_i}\big(h|\hat{H}^*(g,R_{b,m})\big)\overline{\operatorname{Tr}\big(g^kh|R_{b,m}\big)} \\
 &=\frac{1}{N}\Bigg(n_i+\sum_{|g^b|=p_i} \frac{-1}{\varphi (p_i)}\zeta_N^{-kb}+\sum_{ |g^b|=p_i^2} \frac{-1}{\varphi \big(p_i^2\big)}\zeta_N^{-kb}+\dots+\sum_{|g^b|=p_i^{n_i}} \frac{-1}{\varphi \big(p_i^{n_i}\big)}\zeta_N^{-kb}\Bigg).
\end{align*}

Let $\big(k,p_i^{n_i}\big)=p_i^l$. If $s>l$, then $\sum_{|g^b|=p_i^s}\zeta_N^{-kb}$ is a sum of primitive $p_i^{s-l}$-th roots of unity and equals an integer multiple of the M\"obius function $\mu\big(p_i^{s-l}\big)$. If $s\leq l$, then we have
\[
\sum_{|g^b|=p_i^s}\zeta_N^{-kb}=\varphi(p_i^s).
\]
Hence, when $\big(k,p_i^{n_i}\big)=p_i^{n_i}$, we have 
\[ a_{k,p_i}=\frac{1}{N}(n_i-1-\dots-1)=0.\]
When $\big(k,p_i^{n_i}\big)=p_i^{l}$, because $\big(N/p_i^{n_i}\big)=\sum_{d|(N/p_i^{n_i})}\varphi\big(N/p_i^{n_i}d\big)$, we have
\begin{align*}
a_{k,p_i}&=\frac{1}{N}\Bigg(n_i-1-\dots-1+\frac{p_i^l}{\varphi \big(p_i^{l+1}\big)}\Bigg)\\
 &=\frac{(n_i-l)\varphi \big(p_i^{l+1}\big)+p_i^l}{N\varphi \big(p_i^{l+1}\big)} =\frac{p_i^{l+1}(n_i-l)-p_i^l(n_i-l-1)}{p_i^{l+1}\varphi\big(p_i^{n_i}\big)\big(N/p_i^{n_i}\big)}\\
 &=\frac{1}{\sum_{(p_i,d)=1,d|N}\varphi (N/d)}\bigg(n_i-l-\frac{n_i-l-1}{p_i}\bigg).\tag*{\qed}
\end{align*}\renewcommand{\qed}{}
\end{proof}
